\chapter{Reconstruction and execution procedures}
\label{ch:execution}

\section{Ingest, reconcile and construct histories}

The historical procedure starts with a filing manifest. It records documents,
reporting dates, publication dates, revisions and expected statements and
notes. Tagged filings are parsed through an XBRL adapter; PDF and HTML sources
are segmented into text, tables, captions and references. The original bytes
and extraction output remain immutable. Re-running an improved extractor
creates a new extraction version with links to the same source.

For each table, the importer first resolves headers, units, signs, periods and
dimensions. It then proposes concept mappings. A deterministic validator checks
subtotal arithmetic, duplicate facts, comparative consistency and cross-note
references. Every proposed repair creates a new reading. The system may suggest
that an OCR value is 10 rather than 70, but only a reviewed correction or an
explicit alternative allows that reading into a validated model.

Reconciliation proceeds from local totals to movement schedules. A debt bridge
may contain opening debt, issuance, repayment, acquired debt, interest
accretion, currency effects, fair-value changes and closing debt. Every
component receives a movement type. An unexplained residual is stored as a
residual with a source and period; it is not automatically classified as
repayment or issuance. A model can introduce an unknown decomposition of that
residual, subject to evidence and economic constraints.

The reconstruction engine then performs the following finite procedure within
a declared model scope:
\begin{enumerate}
\item Build shared variables for observations, opening states, contract terms
and supported movement amounts. Attach original rounding domains.
\item Add identities, roll-forwards, account restrictions, contractual
relationships and cross-period links. Add each semantic reading as a condition
on the same joint model.
\item Propagate types, units, intervals and simple equalities. Reject context
mismatches before invoking a numerical solver.
\item Ask the solver for feasibility. On infeasibility, return a sufficient
conflicting subset of constraints, its source spans and model assumptions.
A solver's conflict core need not be the smallest possible explanation.
\item Retain finite alternatives and continuous solution families symbolically.
Materialise example histories only when needed for review or execution.
\item Check whether each requested simulation has sufficient economic state.
Missing maturities, reset dates or cohort composition become declared
assumptions or explicit unresolved inputs.
\end{enumerate}

Unrestricted journal enumeration would be unbounded: arbitrary offsetting
transactions can be inserted without changing reports. The model therefore
declares allowed event types, time resolution, relevant accounts, lot/cohort
resolution and limits on latent contract instances. These limits describe the
represented family; they are not evidence that all real possibilities were
searched. The compiler canonicalises syntactic duplicates and equivalences
established by its supported rewrite rules. It does not decide arbitrary
economic equivalence. Probability measures attach to declared economic
hypotheses or parameters, never to a count of encodings; adding a duplicate
representation therefore cannot create extra probability mass or coverage.

The first implementation uses symbolic linear bridges and finite structural
choices. Nonlinear economic models can be evaluated on feasible states through
plugins. A solver timeout preserves the unresolved family and returns
\code{SolverIncomplete}; it does not prune the family as impossible. Search
budgets belong to run configuration and are recorded with the results.

\section{Condition a run without rewriting history}

A run selects an evidence view and cutoff, a historical family, a forecast
starting date, scenario definitions, accounting packages, management policies,
output queries and numerical settings. A run may use a concrete history,
an exact symbolic family, a finite set of conditional histories, or weighted
draws from a declared distribution. The output records which mode was used.

Suppose no debt maturity is disclosed. An analyst may run a conditional case
that assumes all debt matures after the forecast horizon. The run is valid
under that assumption even though the historical representation remains
incomplete. The interface must display the assumption next to the result and
provide sensitivity cases when the assumption is material. It must not label
the assumed maturity as a discovered historical fact.

Management behaviour normally uses only information available at the decision
date. A perfect-foresight stress case is a separate, explicitly labelled
policy. Scenarios should share the same underlying innovations when estimating
paired differences where appropriate, with stable stream identifiers so
branch order does not change the draws. Common random numbers improve the
comparability of paired simulations; they do not imply that every branch has
the same economic event counts or shock distribution.

\section{Schedule and commit events}

Each event carries an economic occurrence time, accrual interval when relevant,
settlement time, entity, contract or cohort identifier, currency, inputs and
provenance. Its ordering key consists of date, declared phase, dependency edges
and a deterministic tie identifier. The reference phases are: opening and
scheduled contractual changes; market observations and rate resets; business
events and decisions; accrual and measurement; settlement; period closing and
reporting. A contract can require a more specific dependency. These phases are
an engineering starting point, not a universal description of market timing.

Two events may share a tie order only if their order is either specified or
irrelevant to the result. If payment before remeasurement and payment after
remeasurement differ, an unresolved ordering is a model ambiguity or an error.
The engine does not conceal it behind alphabetical event IDs.

\begin{lstlisting}
execute_event(event, state, run_context):
    verify prerequisites and information availability
    resolve applicable measurement and recognition rules
    compute proposed economic changes and measurements
    compute proposed journal batch(es)
    validate types, contexts, balance, constraints and idempotency
    atomically commit state, journals and explanation edges
    return new state, emitted dependent events, diagnostics
\end{lstlisting}

The atomic unit includes the economic state and its accounting effects.
If posting fails, neither state changes. A repeated event ID with identical
inputs returns the existing result; the same ID with different inputs is a
conflict. Consolidation runs after entity books reach the required date and
stores its adjustments separately. A failure in one scenario branch cannot
mutate another branch's state.

\section{Funding, tax and other simultaneous relations}

Some financial models have circular dependencies. Interest affects taxable
profit; tax and investment affect funding needs; funding affects interest.
The reference executor first tries to resolve these through explicit timing.
For example, beginning-of-period debt can determine current interest, with
new borrowing affecting the next accrual period. That is acceptable only when
the model's timing assumption is stated and appropriate.

Where the economic model genuinely requires simultaneous decisions, the
dependency graph identifies the strongly connected component and delegates it
to a declared equation or optimisation solver. The solver must return a
solution, residuals, admissibility checks and uniqueness or multiplicity
status. A contraction iteration requires a justified convergence condition;
an arbitrary fixed iteration count is not a mathematical solution. Multiple
admissible equilibria become alternatives when material. No solution produces
a model outcome such as financing failure, or an explicit numerical failure,
depending on the cause.

Funding policy must include available facilities, limits, covenants, collateral
and allowed capital actions. The program cannot issue unlimited debt to keep
cash positive. An inability to meet payments may be a legitimate default
outcome. Default ends or changes contractual and accounting processes under
the relevant event rules; it is not quietly discarded as a bad simulation draw.

\section{Worked example: identical history, different future journals}

Return to the company in Chapter~\ref{ch:purpose}. Immediately after the
historical interest payment and closing, cash is 95, debt is 100 and equity
is $-5$. There are no other assets, liabilities or transactions. Next year's
expense is 5 in the all-fixed branch and 7 in the all-floating branch under
the rate shock. The generated entries and closing statements are:

\begin{center}
\begin{tabular}{lrr}
\toprule Output & All fixed & All floating\\\midrule
Debit interest expense / credit cash &5&7\\
Closing cash &90&88\\
Closing debt &100&100\\
Current-year profit & $-5$ & $-7$\\
Closing equity after profit transfer &$-10$&$-12$\\
Cash change &$-5$&$-7$\\\bottomrule
\end{tabular}
\end{center}
The balance-sheet identity holds in each branch: $90=100-10$ and
$88=100-12$. The demonstration reporting profile classifies interest paid as
operating cash flow. A different permitted reporting profile changes that
classification, while preserving the cash movement and source event. This
example proves the intended distinction and its arithmetic; it does not
estimate the probability that either debt structure describes a real company.

\section{Worked example: a bond has several relevant amounts}

A two-period zero-coupon bond raises cash of 100 and promises 110.25 at maturity.
Assume a 5\% effective rate per period, no transaction costs, no default, and
amortised-cost treatment under a demonstration rule. The effective-interest
schedule follows directly from $100(1.05)^2=110.25$:
\begin{center}
\begin{tabular}{lrrr}
\toprule Date & Opening liability & Interest accretion & Closing liability\\\midrule
Issue &0&0&100\\
End of period 1 &100&5&105\\
End of period 2, before payment &105&5.25&110.25\\
Redemption &110.25&0&0\\\bottomrule
\end{tabular}
\end{center}
Issuance debits cash and credits the liability by 100. Each accretion debits
interest expense and credits the liability. Redemption debits the liability
and credits cash by 110.25. Legal face amount, carrying amount and market
fair value remain separate fields. A market-rate shock may change fair value
without changing this contractual payment schedule or this demonstration
amortised-cost path.

Cash-flow reporting uses the actual payment event and its economic components;
the annual accretion is non-cash. The applicable reporting package decides how
to classify the maturity payment, rather than inferring it solely from the
liability debit. Fees, modifications, early redemption, embedded options,
hedges and foreign currency require additional event and measurement rules.

Ordinary equity issuance similarly separates cash, share quantity and legal
capital. Ten new shares at 12, with par value 1, produce cash of 120, share
capital of 10 and premium of 110 under a simple par-value profile. Shares
outstanding increase by ten units. Direct issuance costs of 3 reduce the
specified equity account and cash under the chosen illustrative treatment.
The accounting rule and the cost's eligibility must be explicit; this example
does not prescribe every jurisdiction's treatment.

\section{Reporting and explanations}

The reporting engine produces the balance sheet, income statement, comprehensive
income, cash-flow statement, changes in equity and note schedules when the
necessary modules are present. Each output carries its basis, perimeter,
period, scenario and interpretation. A comparative table cannot combine an
original prior-year figure and a restated current-year comparative without
displaying that choice.

An explanation is a dependency graph with typed edges: extracted from,
interpreted as, assumed by, constrained by, generated by, measured by,
posted by, aggregated into and supersedes. A cash-flow line must lead back
to payment events; a fair-value movement must lead back to valuation inputs;
a consolidated amount must show eliminations. A provenance path establishes
how a result was obtained. It does not by itself establish that the underlying
economic model is correct.
